% preamble.tex — Shared preamble for the textbook.
% Stack: ctexbook + xelatex + biblatex/biber
% Designed for Chinese prose + English paper titles + math.

\documentclass[a4paper,12pt,openany]{ctexbook}

% ─── Encoding & Fonts ────────────────────────────────────
\usepackage{fontspec}
% Use system fonts for CJK and Latin
% On macOS: Songti SC / STSong for Chinese; falls back to system default
% On Linux: install fonts-noto-cjk
\setCJKmainfont[BoldFont=STHeiti,ItalicFont=STKaiti]{STSong}
\setCJKsansfont{STHeiti}
\setCJKmonofont{STFangsong}
\setmainfont{Times New Roman}
\setsansfont{Arial}
\setmonofont{Courier New}

% ─── Math ────────────────────────────────────────────────
\usepackage{amsmath,amssymb,amsthm,mathtools}
\usepackage{bm}           % bold math
\usepackage{nicefrac}     % compact 1/2 etc.

% ─── Theorem Environments ────────────────────────────────
\newtheorem{theorem}{定理}[chapter]
\newtheorem{lemma}[theorem]{引理}
\newtheorem{corollary}[theorem]{推论}
\newtheorem{proposition}[theorem]{命题}
\theoremstyle{definition}
\newtheorem{definition}[theorem]{定义}
\newtheorem{example}[theorem]{例子}
\newtheorem{exercise}{练习}[chapter]
\theoremstyle{remark}
\newtheorem*{remark}{注}

% ─── Bibliography ────────────────────────────────────────
\usepackage[
  backend=biber,
  style=authoryear-comp,
  sorting=nyt,
  maxbibnames=99,
  maxcitenames=2,
  uniquelist=false,
  uniquename=false,
  natbib=true,
]{biblatex}
\addbibresource{bibliography/library.bib}

% ─── Layout ──────────────────────────────────────────────
\usepackage[
  a4paper,
  top=2.5cm,bottom=2.5cm,
  left=3cm,right=3cm,
]{geometry}
\usepackage{setspace}
\onehalfspacing

% ─── Figures & Tables ────────────────────────────────────
\usepackage{graphicx}
\graphicspath{{figures/}}
\usepackage{booktabs}
\usepackage{multirow}
\usepackage{subcaption}
\usepackage{float}

% ─── Code & Algorithms ───────────────────────────────────
\usepackage{listings}
\lstset{
  basicstyle=\ttfamily\small,
  breaklines=true,
  frame=single,
  language=Python,
}
\usepackage{algorithm}
\usepackage{algorithmicx}
\usepackage{algpseudocode}

% ─── Hyperlinks ──────────────────────────────────────────
\usepackage[colorlinks=true,linkcolor=blue!60!black,citecolor=teal,urlcolor=teal]{hyperref}
\usepackage[capitalise,noabbrev]{cleveref}

% ─── Custom Notation ─────────────────────────────────────
\newcommand{\R}{\mathbb{R}}
\newcommand{\E}{\mathbb{E}}
\newcommand{\inner}[2]{\langle #1,\, #2 \rangle}
\newcommand{\norm}[1]{\left\| #1 \right\|}
\newcommand{\mb}[1]{\boldsymbol{#1}}
\newcommand{\vx}{\boldsymbol{x}}
\newcommand{\vy}{\boldsymbol{y}}
\newcommand{\vq}{\boldsymbol{q}}
\newcommand{\vk}{\boldsymbol{k}}
\newcommand{\vv}{\boldsymbol{v}}
\newcommand{\mW}{\boldsymbol{W}}
\newcommand{\mS}{\boldsymbol{S}}
% TTT-specific
\newcommand{\tttloss}{\mathcal{L}_{\mathrm{TTT}}}
\newcommand{\innerloop}{\mathcal{L}_{\mathrm{inner}}}
\newcommand{\outerloop}{\mathcal{L}_{\mathrm{outer}}}

% ─── Custom Boxes ────────────────────────────────────────
\usepackage{tcolorbox}
\tcbuselibrary{breakable,skins}
\newtcolorbox{toybox}[1][]{
  title={\textbf{玩具问题}},
  colback=blue!3,colframe=blue!40!black,
  fonttitle=\bfseries, breakable, #1
}
\newtcolorbox{misconceptionbox}[1][]{
  title={\textbf{常见误解}},
  colback=orange!5,colframe=orange!60!black,
  fonttitle=\bfseries, breakable, #1
}
\newtcolorbox{neurosciencebox}[1][]{
  title={\textbf{神经科学直觉}},
  colback=purple!5,colframe=purple!50!black,
  fonttitle=\bfseries, breakable, #1
}
\newtcolorbox{analogyboundarybox}[1][]{
  title={\textbf{类比边界与失效条件}},
  colback=red!3,colframe=red!60!black,
  fonttitle=\bfseries, breakable, #1
}
\newtcolorbox{labbox}[1][]{
  title={\textbf{动手实验}},
  colback=green!3,colframe=green!40!black,
  fonttitle=\bfseries, breakable, #1
}
\newtcolorbox{viewpointbox}[1][]{
  title={\textbf{视角（非定论）}},
  colback=purple!3,colframe=purple!40!black,
  fonttitle=\bfseries, breakable, #1
}
\newtcolorbox{insightbox}[1][]{
  title={\textbf{核心洞见}},
  colback=teal!4,colframe=teal!50!black,
  fonttitle=\bfseries, breakable, #1
}
\newtcolorbox{openproblembox}[1][]{
  title={\textbf{开放问题}},
  colback=gray!4,colframe=gray!50!black,
  fonttitle=\bfseries, breakable, #1
}
% Alias: \insight used as environment name in chapters
\newenvironment{insight}[1][]{\begin{insightbox}[#1]}{\end{insightbox}}
\newenvironment{openproblem}[1][]{\begin{openproblembox}[#1]}{\end{openproblembox}}

% ─── Chapter/Section Formatting ──────────────────────────
\usepackage{titlesec}
% (ctexbook handles most CJK chapter formatting automatically)

% ─── Other Utilities ─────────────────────────────────────
\usepackage{enumitem}
\usepackage{microtype}
\usepackage{xcolor}
\usepackage{epigraph}
\setlength{\epigraphwidth}{0.65\textwidth}
% Math macros for chapter 3
\newcommand{\calT}{\mathcal{T}}
\newcommand{\calL}{\mathcal{L}}
\newcommand{\calS}{\mathcal{S}}
\newcommand{\calQ}{\mathcal{Q}}
